\documentclass{ctexart}

% 数学相关宏包
\usepackage{amsmath,mathtools,amsthm,amssymb}
\usepackage{bbm}
\usepackage{mathrsfs}
\usepackage{centernot}

% 图形和绘图
\usepackage{graphicx,float}
\usepackage{tikz,tikz-cd}
\usepackage{pgfplots}
\usepackage[framemethod=TikZ]{mdframed}
\usetikzlibrary{decorations.pathmorphing,positioning,arrows,automata,calc,shapes,patterns}
\pgfplotsset{compat=1.18}
\usepgfplotslibrary{statistics}

% 表格和列表
\usepackage{booktabs,multirow,colortbl}
\usepackage{enumitem}
\setlist[itemize,enumerate,description]{noitemsep}
\setlistdepth{9}

% 算法
\usepackage{algorithm,algorithmic}

% 字体和颜色
\usepackage{xcolor,listings}

% 页面和链接
\usepackage{fancyhdr,caption}
\usepackage{hyperref}
\usepackage{cleveref}

% 工具宏包
\usepackage{etoolbox,letltxmacro,docmute}
\usepackage{gensymb}

% 自定义设置
\renewcommand{\arraystretch}{1.5}
\let\originalmathbb\mathbb
\renewcommand{\mathbb}[1]{%
    \ifstrequal{#1}{1}{\mathbbm{#1}}{\originalmathbb{#1}}%
}

% 颜色定义
\definecolor{lightblue}{RGB}{173,216,230}
\definecolor{lightgreen}{RGB}{144,238,144}
\definecolor{lightyellow}{RGB}{255,255,224}
\definecolor{lightgray}{RGB}{211,211,211}

% 超链接设置
\hypersetup{
    colorlinks=true,
    linkcolor=red,
    filecolor=magenta,
    urlcolor=cyan,
    pdftitle={数学笔记},
    pdfauthor={乐绎华},
    pdfsubject={数学笔记},
    pdfkeywords={数学},
    bookmarksnumbered=true,
    bookmarksopen=true,
    bookmarksopenlevel=6,
    pdfstartview=FitH,
    pdfpagemode=UseOutlines,
    pdfborder={0 0 0}
}

% 定理环境
\theoremstyle{plain}
\newtheorem{theorem}{定理}[section]
\newtheorem{lemma}[theorem]{引理}
\newtheorem{corollary}[theorem]{推论}
\newtheorem{proposition}[theorem]{命题}

\theoremstyle{definition}
\newtheorem{definition}[theorem]{定义}
\newtheorem{example}[theorem]{例}
\newtheorem{exercise}[theorem]{练习}
\newtheorem{problem}[theorem]{题目}

\theoremstyle{remark}
\newtheorem{remark}[theorem]{备注}
\newtheorem{note}[theorem]{注}
\newtheorem{observation}[theorem]{观察}

% 解答环境
\newtheorem*{solution}{解}
\newtheorem*{proof*}{证明}

% 图片路径
\graphicspath{{F:/iCloudDrive/iCloud~md~obsidian/2025-summer/figures/}}

\author{乐绎华}
\title{数学笔记}
\date{\today}

\begin{document}

\maketitle
\tableofcontents

See loring W. tu .

\begin{figure}[H]
\centering
\includegraphics[width=\textwidth]{de-rham-cohomology-2025073100.png}
% \caption{}
\label{}
\end{figure}
\begin{figure}[H]
\centering
\includegraphics[width=\textwidth]{1-de-rham-cohomology-2025073100.png}
% \caption{}
\label{}
\end{figure}
\begin{figure}[H]
\centering
\includegraphics[width=\textwidth]{2-de-rham-cohomology-2025073100.png}
% \caption{}
\label{}
\end{figure}

The de Rham cohomology $H^{*}(M)$ is an anticommutative graded ring, with wedge as the multiplication; anticommutative means :
\[
a\cdot b=(-1)^{k\ell}b\cdot a
\]
where $a\in A^{k}, b\in A^{\ell}$ for the graded ring $A=\oplus_{k=0}^{\infty}A^{k}$; in graded ring the ring multiplication sends $A^{k}\times A^{\ell}$ to $A^{k\times \ell}$.

Suppose $F:N\to M$ is a $C^{\infty}$ mapping of manifolds; because $F^{*}(\omega \wedge \tau)=F^{*}\omega \wedge F^{*}\tau$ for differential forms $\omega$ and $\tau$ on $M$, the pullback map $F^{*}:H^{*}(M)\to H^{*}(N)$ is a ring homomorphism. If $F$ is diffeomorphism, then $F^{*}$ is isomorphism.

To sum up, de Rham cohomology gives a contravariant functor from the category of $C^{\infty}$ manifolds to the category of anticommutative graded rings. In this way the de Rham cohomology becomes a powerful diffeomorphism invariant of $C^{\infty}$ manifolds.

\begin{figure}[H]
\centering
\includegraphics[width=\textwidth]{de-rham-cohomology-2025073112.png}
% \caption{}
\label{}
\end{figure}


\end{document}